Dans tout le chapitre, on prend $I$ un intervalle réel tel que $\mathring{I} \neq \emptyset$.

On prendra aussi $\K = \R$ ou $\C$, et $(E,N)$ un $\K$-espace vectoriel normé
 de dimension finie.

\section{Dérivation}
\subsection{Dérivée en un point}

\begin{definition}{Dérivable en un point}{}
    Soit $a \in I$ et $f: I \to E$. On dit que $f$ est dérivable en $a$ ssi 
    \[\fonction{\Delta_a}{I \setminus \{a\}}{E}{t}{\frac{f(t) - f(a)}{t-a}}\]
    a une limite $\ell \in E$ en $a$. Cette limite est appelée dérivée (ou vecteur dérivé)
    de $f$ en $a$ et se note $f'(a)$. \\

    On a alors : $\displaystyle \frac{f(a + h ) - f(a)}{h} \xrightarrow[h \to 0]{N} f'(a)$.
\end{definition}

\begin{proposition}{Caractérisation de la dérivabilité}{}
    Soit $a \in I$ et $f: I \to E$. Alors $f$ est dérivable en $a$ ssi il existe $V \in E$
    et $\varepsilon : I \to E$, avec $\varepsilon(t) \xrightarrow[t \to a]{} 0$ telle que
    \[ \forall t \in I,~f(t) = f(a) + (t-a)V + (t-a)\varepsilon(t)\]

    appelé développement limité de $f$ en $a$ à l'ordre 1, et alors $V = f'(a)$.
\end{proposition}

\begin{proof}
    $\Rightarrow$ Si $f$ est dérivable en $a$, posons $V = f'(a)$ et $\varepsilon(t) = \frac{f(t) - f(a)}{t - a } - f'(a)$ si 
    $t \neq a$ et $\varepsilon(a) = 0$. 

    Alors $\varepsilon$ vérifie l'égalité et $\varepsilon$ continue en $a$ car 
    $\displaystyle \frac{f(t) - f(a)}{t-a} \xrightarrow[t \to a]{} f'(a)$, c'est à dire 
    $\varepsilon(t) \xrightarrow[t \to a]{} 0$.

    $\Leftarrow$ Si $\varepsilon$ et $V$ existent, $\displaystyle \frac{f(t) - f(a)}{t-a} = V + \varepsilon(t) \limit t a V$
    Donc $f$ dérivable en $a$ et $f'(a) = V$.
\end{proof}

\begin{remarque}
Avoir un DL2 n'implique pas d'être deux fois dérivable.
Par exemple :
\[\fonction{f}{\R}{\R}{x}{
    \begin{cases} x^3\sin\left(\frac 1 {x^2}\right)~\text{si}~x \not = 0 \\
    0~\text{sinon}
    \end{cases}}\]
a un DL2 mais n'est pas dérivable deux fois.
\end{remarque}
\begin{corollaire}{Dérivabilité implique continuité}{}
    Si $f$ est dérivable en $a$, alors $f$ est continue en $a$.
\end{corollaire}

\begin{proof}
    En ce cas $f(t) = f(a) + (t-a)f'(a) + (t-a)\varepsilon(t) \limit{t}{a} f(a)$.
\end{proof}


\begin{definition}{Interprétation cinématique}{}
    Soit $f : I \to E$ une courbe paramétrée. 

    \begin{itemize}
        \item Soit $t$ un instant et $f(t)$ la position d'un point matériel à l'instant 
        $t$, alors $f'(t)$ est le vecteur vitesse à l'instant $t$. 

        \item Si $f'(a) \neq 0$, $T = \{ f(a) + hf'(a), \, h \in \R\}$ est la tangeante à la courbe à l'instant 
        $a$. 

        \item $C = \{ f(t), \, t \in I\}$ est le support de la courbe paramétrée par $f$.  
    \end{itemize}
\end{definition}

\begin{proposition}{Lien avec une base}{}
    Soit $B = (e_1, \ldots, e_n)$ une base de $E$, et soit 
    \[\fonction{f}{I}{E}{t}{\sum_{i=1}^n f_i(t)e_i},~\text{où}~f_i : I \to \R\]

    Soit $a \in I$. $f$ est dérivable en $a$ ssi $\forall i \in \ninter{1}{n}~f_i$
    est dérivable en $a$, et alors :
    \[f'(a) = \sum_{i=1}^n f_i'(a)e_i\]
\end{proposition}

\begin{proof}
    $\Delta_a(t) = \displaystyle \frac{f(t) - f(a)}{t-a} = \sum_{i=1}^n \frac{f_i(t) - f_i(a)}{t-a}e_i$ a 
    une limite quand $t \to a $ ssi toutes ses coordonnées ont une limite.
\end{proof}

\begin{exemple}
    $\fonction f \R {\R^3} t {(\cos t, \sin t, t)}$ représente une hélice.

    $\fonction g \R {\R^2} t {(\ch t, \sh t)}$ est une hyperbole.

    La lemniscate de Bernoulli est aussi une courbe paramétrée, plus précisément par $\rho = \sqrt{\cos {2\theta}}$.
\end{exemple}

\begin{figure}[htbp]
    \centering
        \begin{minipage}[t]{0.32\textwidth}
        \centering
        \begin{tikzpicture}
            \begin{axis}[
                width=\linewidth,
                height=4.8cm,
                view={65}{30},
                axis lines=center,
                axis line style={-stealth, black!50, line width=0.4pt},
                ticks=none,
                xlabel={$x$}, ylabel={$y$}, zlabel={$z$},
                label style={font=\footnotesize},
                xmin=-1.3, xmax=1.3,
                ymin=-1.3, ymax=1.3,
                zmin=0,    zmax=14,
                enlargelimits=0.1
            ]
            \addplot3[
                domain=0:4*pi,
                samples=160,
                samples y=0,
                vorange,
                line width=1.2pt
            ] (
                {cos(deg(x))},
                {sin(deg(x))},
                {x}
            );
            \end{axis}
        \end{tikzpicture}
        \par\vspace{2pt}
        {\footnotesize\textbf{Hélice}\par}
        {\scriptsize $t \mapsto (\cos t, \sin t, t)$\par}
        \end{minipage}\hfill
        \begin{minipage}[t]{0.32\textwidth}
            \centering
            \begin{tikzpicture}
                \begin{axis}[
                    width=\linewidth,
                    height=4.8cm,
                    axis lines=middle,
                    axis line style={-stealth, black!50, line width=0.4pt},
                    ticks=none,
                    xlabel={$x$}, ylabel={$y$},
                    label style={font=\footnotesize},
                    xmin=-0.5, xmax=3.2,
                    ymin=-2.5, ymax=2.5,
                    axis equal image
                ]
                \addplot[densely dashed, black!35, line width=0.5pt, domain=0:2.5] {x};
                \addplot[densely dashed, black!35, line width=0.5pt, domain=0:2.5] {-x};
                \addplot[
                    domain=-1.5:1.5,
                    samples=80,
                    vorange,
                    line width=1.2pt
                ] (
                {cosh(x)},
                {sinh(x)}
                );
                \end{axis}
            \end{tikzpicture}
            \par\vspace{2pt}
            {\footnotesize\textbf{Hyperbole}\par}
            {\scriptsize $t \mapsto (\ch t, \sh t)$\par}
            \end{minipage}\hfill
            \begin{minipage}[t]{0.32\textwidth}
            \centering
            \begin{tikzpicture}
                \begin{axis}[
                    width=\linewidth,
                    height=4.8cm,
                    axis lines=middle,
                    axis line style={-stealth, black!50, line width=0.4pt},
                    ticks=none,
                    xlabel={$x$}, ylabel={$y$},
                    label style={font=\footnotesize},
                    xmin=-1.25, xmax=1.25,
                    ymin=-0.65, ymax=0.65,
                    axis equal image
                ]
                \addplot[
                    domain=0:360,
                    samples=160,
                    vorange,
                    line width=1.2pt
                ] (
                {cos(x) / (1 + sin(x)^2)},
                {sin(x) * cos(x) / (1 + sin(x)^2)}
                );
                \end{axis}
            \end{tikzpicture}
            \par\vspace{2pt}
            {\footnotesize\textbf{Lemniscate de Bernoulli}\par}
            {\scriptsize $\rho^2 = \cos(2\theta)$\par}
        \end{minipage}
\end{figure}

\begin{definition}{Dérivabilité à droite et à gauche}{}
    Soit $f : I \to E$ et $a \in \mathring I$. 

    \avspace On définit $\fonction{f_g}{I \cap ]-\infty, a]}{E}{t}{f(t)}$ et
    $\fonction{f_d}{I \cap [a, +\infty[}{E}{t}{f(t)}$ \\

    Comme $a \in \mathring I$, $I \cap ]-\infty, a]$ et $I \cap [a, +\infty[$ 
    sont d'intérieur non vide. \\

    On dit que $f$ est dérivable à gauche (resp. à droite) en $a$ ssi $f_g$ 
    (resp. $f_d$) est dérivable en $a$.
    On définit alors le vecteur dérivé à gauche (resp. à droite) de $f$ en $a$ par :
    \[f'_g(a) = \lim_{t \to a} \frac{f_g(t) - f_g(a)}{t-a} = \lim_{t \to a^-} \frac{f(t) - f(a)}{t-a}~\text{(resp. à droite par}~f'_d(a) = \lim_{t \to a^+} \frac{f(t) - f(a)}{t-a}\text{)}\]                                                    
\end{definition}

\begin{proposition}{Dérivabilité en un point avec dérivées à gauche et à droite}{}
    $f$ est dérivable en $a$ ssi elle est dérivable à gauche et à droite en $a$ et
    $f'_g(a) = f'_d(a)$. 
\end{proposition}

\begin{proof}
    Coordonnée par coordonnée.
\end{proof}

\begin{definition}{Point anguleux - LHP}{}
    Soit $f : I \to E$, et $a \in \mathring{I}$. 
    On suppose que $f$ est dérivable à gauche et à droite en $a$ et que 
    $f'_g(a) \neq f'_d(a)$.

    On dit alors que la courbe paramétrée par $f$ possède en $a$ un point anguleux.
\end{definition}

\cleanmarginenote{
  \begin{tikzpicture}[scale=0.9, >=stealth]
    \draw[->, black!35, line width=0.4pt] (-0.3, 0) -- (2.3, 0) node[right, font=\tiny] {$x$};
    \draw[->, black!35, line width=0.4pt] (0, -2.1) -- (0, 1.3) node[above, font=\tiny] {$y$};

    \draw[vorange, line width=1.3pt] (0, -1.8) -- (0, 0) -- (1.8, 0);

    \draw[->, dangerred, line width=1.1pt] (0, 0) -- (0, 1) 
      node[left=1pt, font=\scriptsize, text=dangerred] {$f'_g(0)$};
    \draw[->, forestgreen!85!black, line width=1.1pt] (0, 0) -- (1, 0) 
      node[above=1pt, font=\scriptsize, text=forestgreen!85!black] {$f'_d(0)$};

    \fill[black] (0, 0) circle (1.2pt) node[below left=1pt, font=\tiny] {$O$};
  \end{tikzpicture}%
}

\begin{exemple}
    $\fonction f \R {\R^2} t {\begin{cases} (t, 0)~\text{si}~x > 0\\(0,t)~\text{sinon}\end{cases}}$ admet un point anguleux en $0$.
\end{exemple}

\subsection{Dérivabilité sur un intervalle}

\begin{definition}{Dérivable sur un intervalle}{}
    Soit $f : I \to E$. On dit que $f$ est dérivable sur $I$ ssi elle l'est en tout $a \in I$. 
    On définit alors sa dérivée :
    \[\fonction{f'}{I}{E}{t}{f'(t)}\] 

    L'ensemble des fonctions dérivables de $I$ dans $E$ est noté $\mathscr D(I, E)$.
\end{definition}

\begin{definition}{Fonction de classe $\mathscr C^1$}{}
    On dit que $f : I \to E$ est de classe $\mathscr C^1$ sur $I$ ssi 
    elle y est dérivable et si $f'$ est continue.
    L'ensemble de ces fonctions est noté $\mathscr C^1(I, E)$.
\end{definition}

\begin{proposition}{Dérivabilité et classe $\mathscr C^1$ avec une base }{}
    Soit $B = (e_1, \ldots, e_n)$ une base de $E$, et soit
    $\displaystyle \fonction{f}{I}{E}{t}{\sum_{i=1}^n f_i(t)e_i}$, où $f_i : I \to \K$. \\

    Alors $f$ est dérivable (resp. de classe $\mathscr C^1$) sur $I$ ssi
    $\forall i \in \ninter{1}{n}$, $f_i$ est dérivable (resp. de classe $\mathscr C^1$).
\end{proposition}

\begin{exemple}
    $\fonction{f}{\R}{\R^2}{t}{(x(t), y(t))}$ avec $x(t) = 2 \cos(t)$ et $y(t) = \sin(t)$.

    Alors $f$ est $\mathscr C^1$ sur $\R$ car $x$ et $y$ le sont.
\end{exemple}

\marginenote{Un peu de Latin}{\emph{Nihil novi sub sole} - rien de nouveau pour le moment.}
\begin{proposition}{Structure fonctions dérivables et $\mathscr C^1$}{}
    $\mathscr D(I, E)$ et $\mathscr C^1(I, E)$ sont des $\K$-espaces vectoriels et :

    \[\fonction{\varphi}{\mathscr C^1(I, E)}{\mathscr C(I, E)}{f}{f'} \text{est linéaire.}\] 
\end{proposition}

\begin{proof}
    Coordonnée par coordonnée.
\end{proof}

\subsection{Composition}

\begin{proposition}{Composition à droite}{}
    Soit $f \in \mathscr C^1(I, E)$ et $\varphi \in \mathscr C^1(J, I)$, où $J$ est un intervalle de $\R$
    tel que $\mathring{J } \neq \emptyset$. 
    Considérons :
    \[\fonction{g = f \circ \varphi}{J}{E}{u}{f(\varphi(u))}\]

    Alors $g \in \mathscr C^1(J, E)$ et $g' = \varphi' f' \circ \varphi$
    
\end{proposition}

\begin{remarque}
    (hors programme)
    \begin{itemize}
        \item La composition par $\varphi$ s'appelle un changement de paramétrage.

        \item Si $\varphi$ est bijective et si $\varphi^{-1} \in \mathscr C^1(I, J)$, c'est à dire si 
        $\varphi '$ ne s'annule pas (on dit alors que $\varphi$ est un $\mathscr C^1$-difféomorphisme). 

        On dit alors que $\varphi$ est un changement de paramétrage admissible.
    \end{itemize}
\end{remarque}

\idee{Coordonnée par coordonnée.}
\begin{proof}
    Soit $B = (e_1, \ldots, e_n)$ une base de $E$. 

    Notons $\displaystyle \forall t \in I, \, f(t) = \sum_{i=1}^n f_i(t)e_i$. 

    Alors $\displaystyle \forall u \in J, \, g(u ) = \sum_{i=1}^n f_i(\varphi(u))e_i$. 

    Donc $g \in \mathscr C^1(J, E)$ et $\displaystyle g'(u) = \sum_{i=1}^n \varphi'(u)f_i'(\varphi(u))e_i = \varphi'(u)f'(\varphi(u))$.
\end{proof}

\begin{theoreme}{Composition à gauche par une application linéaire}{}

    Soit $(F, N)$ un espace vectoriel normé de dimension finie, $f$ dérivable en $a$
    (resp. sur $I$, resp. de classe $\mathscr C^1$ sur $I$) et $u \in \mathscr{L }(E, F)$. \\

    Alors $u \circ f$ est dérivable en $a$ (resp. sur $I$, resp. de classe $\mathscr C^1$ sur $I$) et :
    \[(u \circ f)'(a) = u \circ f'(a)~\text{(resp.} (u \circ f)' = u \circ f'\text{)}\]
\end{theoreme}

\idee{Passer par le DL de $f$, et composer par $u$.}

\begin{proof}
    \uline{En $a$ :}

    Par hypothèse, pour $t \in I$, $f(t) = f(a) + (t-a)f'(a) + (t-a)\varepsilon(t)$

    En composant par $u$, on obtient $u(f(t)) = u(f(a)) + (t-a)u(f'(a)) + (t-a)u(\varepsilon(t))$.

    Et on a bien $u(\varepsilon(t)) \xrightarrow[t \to a]{} 0$, donc $u \circ f$ dérivable en $a$ et $(u \circ f)'(a) = u(f'(a))$.

    Or ceci est vrai en tout point de $I$, donc $(u \circ f)' = u \circ f'$.

    Si $f \in \mathscr C^1(I, E)$, $u$ étant linéaire entre espaces de dimension finie, elle est continue donc $(u \circ f)'$ 
    est continue. 
\end{proof}

\begin{exemple}
    Résolvons le système 
    \[(S) = \begin{cases}
        x'(t) = 2x(t) + y(t) \\
        y'(t) = 4x(t) - y(t)
    \end{cases}\]

    Posons $X(t) = \begin{pmatrix} x(t) \\ y(t) \end{pmatrix}$ et $x, y$ sont $\mathscr C^1$ 
    de $\R$ dans $\R$. De même, $X$ est $\mathscr C^1$ de $\R$ dans $\R^2$. Donc : 

    $$(S) \Leftrightarrow \forall t \in \R, \, X'(t) = \begin{pmatrix} 2 & 1 \\ 4 & -1 \end{pmatrix} X(t)$$

    Posons $P = \begin{pmatrix} 1 & 1 \\ 1 & -4 \end{pmatrix}$, $P$ est inversible car $\det(P) = -5 \neq 0$.

    Posons $\forall t, \, Y(t) = P^{-1}X(t)$. Comme $\fonction u {\R^2}{\R^2}{M}{P^{-1}M}$ est linéaire, 
    $Y$ est $\mathscr C^1$ et $Y'(t) = P^{-1}X'(t)$ car $Y = u \circ X$ donc $Y' = u \circ X'$.

    Donc $(S) \Leftrightarrow \forall t \in \R, \, Y'(t) = P^{-1} \begin{pmatrix} 2 & 1 \\ 4 & -1 \end{pmatrix} P y(t)$

    Posons $Y(t) = \begin{pmatrix} \varphi(t) \\ \psi(t) \end{pmatrix}$ où 
    $\varphi, \psi$ sont $\mathscr C^1$ de $\R$ dans $\R$. Alors : 

    \begin{align*}
        (S) & \Leftrightarrow \forall t \in \R, \, \begin{pmatrix} \varphi'(t)  = 3\varphi(t) \\ \psi'(t) = -2\psi(t) \end{pmatrix}\\
    & \Leftrightarrow \alpha, \beta \in \R \quad \varphi(t) = \alpha e^{3t} \space \psi(t) = \beta e^{-2t}\\
    & \Leftrightarrow \alpha, \beta \forall t, \, X(t) = P \begin{pmatrix} \alpha e^{3t} \\ \beta e^{-2t} \end{pmatrix} = \begin{pmatrix} \alpha e^{3t} + \beta e^{-2t} \\ \alpha e^{3t} - 4\beta e^{-2t} \end{pmatrix}
    \end{align*}
    
\end{exemple}

\marginenote[0.5cm]{Tenseurs - HP}{En définissant $E\otimes F$ le produit \emph{tensoriel} de $E$ et $F$ on pourrait créer des applications linéaires à partir des bilinéaires.}
\begin{theoreme}{Composition à gauche par une application bilinéaire}{}

    Soient $f : I \to E$, $g : I \to F$ des applications et $B : E \times F \to G$ une application bilinéaire où 
    $(G, N'')$ est un espace vectoriel normé de dimension finie.
    Notons :
    \[\fonction{B(f, g)}{I}{G}{t}{B(f(t), g(t))}\]

    Soit $a \in I$. Si $f$ et $g$ sont dérivables en $a$ (resp. sur $I$, resp. de classe $\mathscr C^1$ sur $I$),
    alors $B(f, g)$ aussi et :
    \[(B(f, g))'(a) = B(f'(a), g(a)) + B(f(a), g'(a))~\text{(resp.}~(B(f, g))' = B(f', g) + B(f, g')\text{)}\]
\end{theoreme}

\idee{DL de $f$ et $g$, puis bilinéarité.}

\begin{proof}
    \uline{En $a$ :}

    Pour $t \in I$, $f(t) = f(a) + (t-a)f'(a) + (t-a)\varepsilon_1(t)$ et 
    $g(t) = g(a) + (t-a)g'(a) + (t-a)\varepsilon_2(t)$.

    Et alors par bilinéarité, 
     
    $B(f(t), g(t)) = B(f(a), g(a)) + (t-a)(B(f(a), g'(a)) + B(f'(a), g(a))) + (t-a)\varepsilon_3(t)$

    où $\varepsilon_3(t) = B(f(a), \varepsilon_2(t)) + (t - a)(B(f'(a), g'(a)) + B(f'(a), \varepsilon_2(t)) + B(\varepsilon_1(t), g'(a)) + B(\varepsilon_1(t), \varepsilon_2(t))) + B(\varepsilon_1(t), g(a))$.

    Or $B$ est continue car bilinéaire en dimention finie, donc 
    $\varepsilon_3(t) \xrightarrow[t \to a]{} B(f(a), 0) + B(f'(a), g'(a)) + B(0, g(a)) + 0 (B(f'(a), 0) + B(0, g'(a)) + B(0, 0)) = 0$

    Donc $B(f, g)$ dérivable en $a$ et $(B(f, g))'(a) = B(f'(a), g(a)) + B(f(a), g'(a))$.

    Sur $I$ : ceci est vrai en tout point de $I$, donc $(B(f, g))' = B(f', g) + B(f, g')$.

    Si $f, g \in \mathscr C^1(I, E)$, $B(f, g)'$ est continue car $B$ est bilinéaire entre espaces de dimension finie, donc continue. 
\end{proof}

\begin{proposition}{}{}
    \begin{itemize}
        \item Le produit scalaire de deux applications dérivables est dérivable.
        \item Le produit matriciel de deux applications dérivables est dérivable.
        \item Si $(E, \langle \cdot, \cdot \rangle)$ est préhilbertien (c'est-à-dire euclidien car de dimension finie), et $f \in \mathscr{C }^1(I, E)$, alors $||f||^2 \in \mathscr C^1(I, \R)$ et
        $(||f||^2)' = 2 \langle f, f' \rangle$.
    \end{itemize}
\end{proposition}

\idee{Bilinéarité.}

\begin{proof}
    Un produit scalaire est bilinéaire, et le produit matriciel également, et si 
    $||f||^2 = \langle f, f \rangle$, donc $(||f||^2)' = \langle f', f \rangle + \langle f, f' \rangle = 2 \langle f, f' \rangle$.
\end{proof}

\begin{exemple}
Soit $f: I \to E$ dérivable telle que $\forall t \in I, \|f(t)\| = K > 0$.

Alors $(\|f\|^2)' = 2 \langle f',f \rangle$. Donc $\mathscr C_f \perp C_{f'}$.
\end{exemple}

\begin{figure}[htbp]
\centering
\begin{tikzpicture}[scale=1.2, >= stealth]
    \draw[->, black!35, line width=0.4pt] (-1.2, 0) -- (1.2, 0) node[right, font=\tiny] {$x$};
    \draw[->, black!35, line width=0.4pt] (0, -1.2) -- (0, 1.2) node[above, font=\tiny] {$y$};

    \draw[thick, vorange] (0,0) circle (1);

    \draw[<->, thick, dangerred] (0,0) -- node[midway, above] {$K$} (1,0);

    \draw[->, forestgreen!85!black, line width=1.1pt] (1, 0) -- (1, 1) 
      node[above=1pt, font=\scriptsize, text=forestgreen!85!black] {$f'(0)$};

    \fill[black] (0, 0) circle (1.2pt) node[below left=1pt, font=\tiny] {$O$};
\end{tikzpicture}
\end{figure}

\begin{theoreme}{Généralisation à une application n-linéaire}{}

    Soient $(E_1, N_1), \ldots, (E_n, N_n), (F, N')$ des espaces vectoriels normés de dimension finie, et
    $f_i : I \to E_i$, $M : E_1 \times \cdots \times E_n \to F$ une application $n$-linéaire.
    Notons :
    \[\fonction{M(f_1, \ldots, f_n)}{I}{F}{t}{M(f_1(t), \ldots, f_n(t))}\]

    \avspace Si $f_1, \ldots, f_n$ sont dérivables en $a$ (resp. sur $I$, resp. de classe $\mathscr C^1$ sur $I$),
    alors $M(f_1, \ldots, f_n)$ aussi et :
    $$(M(f_1, \ldots, f_n))'(a) = \sum_{i=1}^n M(f_1(a), \ldots, f_{i-1}(a), f_i'(a), f_{i+1}(a), \ldots, f_n(a))$$
\end{theoreme}

\begin{proof}
    On généralise celle réalisée pour une application bilinéaire via une récurrence.
\end{proof}

\begin{corollaire}{}{}
    Si $f \in \mathscr C^1(I, \mathscr M_n(\K))$ alors $\det \circ f \in \mathscr C^1(I, \K)$.
\end{corollaire}

\begin{exemple}
    Soit $f: t \to \det(X_1(t), X_2(t))$ avec $X_1, X_2 \in \mathscr F(\R, \mathscr M_{2, 1}(\K))$. 

    Alors si $X_1, X_2$ sont dérivables, $f$ aussi et 
    $f' = \det(X_1', X_2) + \det(X_1, X_2')$.
\end{exemple}

\subsection{Application de classe \texorpdfstring{$\mathscr C^k$}{Ck}}

\begin{definition}{Fonction de classe $\mathscr C^k$}{}
    Soit $f : I \to E$ et $k \in \N, \, k \geqslant 2$. On dit que $f$ est de classe 
    $\mathscr C^k$ sur $I$ ssi $f, f', \ldots, f^{(k-1)}$ sont dérivables sur $I$ et $f^{(k)}$ est continue sur $I$. \\

    On dit que $f$ est de classe $\mathscr C^\infty$ sur I ssi $f$ est de classe 
    $\mathscr C^k$ pour tout $k \in \N$.
\end{definition}

\begin{remarque}
    Si $f \in \mathscr C^k(I,E)$ alors :
    $\ds \forall i,j \in \N$ tels que $i + j \leqslant k$,
    $\left(f^{(i)}\right)^{(j)} = \left(f^{(j)}\right)^{(i)} = f^{(i+j)}$.
\end{remarque}

\begin{proposition}{Classe $\mathscr C^k$ avec une base}{}
    Soit $B = (e_1, \ldots, e_n)$ une base de $E$, et :
    \[\fonction{f}{I}{E}{t}{\sum_{i=1}^n f_i(t)e_i},~\text{où}~f_i : I \to \K\]
    

    Alors $f$ est de classe $\mathscr C^k$ sur $I$ ssi $\forall i \in \ninter{1}{n}$, $f_i$ est de classe $\mathscr C^k$. 

    On a alors $\displaystyle \forall j \in \ninter{0}{k}, \, f^{(j)} = \sum_{i=1}^n f_i^{(j)}e_i$.
\end{proposition}

\idee{Récurrence sur $k$.}

\begin{proof}
    Par récurrence sur $k$. 

    \uline{Initialisation :} c'est vrai pour $k = 1$ (vu en sup).

    \uline{Hérédité :} supposons le pour toute fonction de classe $\mathscr C^{k - 1} $. Soit 
    $f \in \mathscr C^k$. Donc $f \in \mathscr C^{k-1}$ ssi $\forall i, \, f_i \in \mathscr C^{k-1}$ 
    et alors $\displaystyle \forall t \in I, \, f^{(k-1)}(t) = \sum_{i=1}^n f_i^{(k-1)}(t)e_i$.

    Et alors en appliquant le cas $k = 1$, on a $f \in \mathscr C^k$ ssi $f^{(k-1)} \in \mathscr C^1$, 
    ssi $\forall i, \, f_i^{(k-1)} \in \mathscr C^1$, ssi $\forall i, \, f_i \in \mathscr C^k$.
\end{proof}

\begin{proposition}{Structure de l'ensemble des fonctions de classe $\mathscr C^k$}{}

    $\mathscr C^k(I, E)$ est un $\K$-espace vectoriel et
    $\fonction{\varphi}{\mathscr C^k(I, E)}{\mathscr C^{k-1}(I, E)}{f}{f'}$ est linéaire.
\end{proposition}

\begin{proof}
    Immédiate en regardant coordonnée par coordonnée.
\end{proof}

\subsection{Rappels : prolongements}

\begin{theoreme}{Limite de la dérivée}{}
    Soit $f : [a, b] \to E$ continue sur $[a, b]$ et dérivable sur $]a, b]$. 

    Si $f'$ a une limite finie $\ell \R$ en $a$, alors $f$ est dérivable en $a$ et $f'(a) = \ell$.
\end{theoreme}

\begin{proof}
    Le faire dans $\R$ puis faire coordonnée par coordonnée.
\end{proof}

\begin{remarque}
    Si $f$ est définie sur $]a,b]$ on peut alors essayer de prolonger $f$ par continuité avant d'appliquer ce théorème.
\end{remarque}

\begin{theoreme}{Prolongement des fonctions de classe $\mathscr C^1$}{}
    Soit $f \in \mathscr C^1(]a, b], \R)$ avec $f$ et $f'$ qui ont des limites finies 
    $\ell_1$ et $\ell_2$ en $a$. 

    Alors on peut prolonger $f$ par continuité en $a$ en posant $f(a) = \ell_1$.

    Et alors $f$ est $\mathscr C^1$ sur $[a, b]$ et $f'(a) = \ell_2$.
\end{theoreme}

\begin{theoreme}{Prolongement des fonctions de classe $\mathscr C^k$ et $\mathscr C^{\infty}$}{}
    Soit $f \in \mathscr C^k(I, \R)$ avec $f,\dots,f^{(k)}$ qui ont des limites finies en $a \in \overline{I} \setminus I$. 

    Alors on peut prolonger $f$ par continuité en $a$.
    Et alors $f$ est $\mathscr C^k$ sur $I \cup \{a\}$ (de même pour $\mathscr C^{\infty}$).
\end{theoreme}

\begin{exemple}[Classique à retenir]
    Soit $\fonction f \R \R x {\begin{cases}
        0 & \text{si } x \leqslant 0\\
        e^{-\frac{1}{x^2}} & \text{sinon}
    \end{cases}}$

    Montrons que $f$ est $\mathscr C^{\infty}$ par prolongement.

    Montrons par récurrence que $\forall x \in \R_+^*, \, f^{(n)}(x) = \frac{P_n(x) }{x^{3n}} e^{-\frac{1}{x^2}}$ où $P_n \in \R [X]$. 

    $f$ est $\mathscr C^\infty$ sur $\R_+^*$. 

    Pour $x > 0$, $f^{(0)}(x) = \frac{P_0(x )}{x^0} e^{-\frac{1}{x^2}}$ avec $P_0 = 1$.

    \uline{Hérédité :} Supposons que $\displaystyle \exists P_n, \, \forall x > 0, \, f^{(n)}(x) = \frac{P_n(x)}{x^{3n}} e^{-\frac{1}{x^2}}$.

    Alors $\displaystyle f^{(n+1)}(x) = \frac{P'_n(x)}{x^{3n}} e^{-\frac{1}{x^2}}  - \frac{3n P_n(x)}{x^{3n + 1}}e^{-\frac{1}{x^2}} - \frac 2 {x^3} \frac{P_n(x)}{x^{3n}}e^{-\frac{1}{x^2}}$

    $f^{(n+1)}(x) = \frac{e^{\frac{-1}{x^2}}}{x^{3n + 3}} \underbrace{\left( x^3 P'_n(x) - 3nx^2 P_n(x) - 2P_n(x) \right)}_{P_{n+1}(x)}$

    On a bien $\forall n, \, f^{(n)}(x) \limit x {0^-} 0 $ et $f^{(n)}(x) \limit x {0^+} 0$ par croissance comparée. 

    Donc par le théorème de prolongement des fonctions de classe $\mathscr C^n$, $f$ est $\mathscr C^n$ sur $\R$ et $\forall n \in \N, \, f^{(n)}(0) = 0$.


    Soit $[a, b] \subset \R$, et $\fonction g \R \R t {f(t-a)f(b-t)}$.

    Alors $g$ est $\mathscr C^\infty$ sur $\R$, $g$ est strictement positive sur $]a, b[$ et nulle ailleurs.

\end{exemple}

\begin{remarque}
    (hors programme)
    
    On dira alors que $g$ est $\mathscr C^{\infty}$ à support compact (plus précisément à support inclus dans un compact).

    L'ensemble de ces fonctions est appelé \emph{espace de Schwartz}, est un espace vectoriel et est noté $S$.

    $\forall f \in \mathscr C^0(\R, \R)$ posons :
    \[\fonction {\varphi_f} S {\R} g {\int_{\R}} f(t)g(t) \dif t\]
    On montrerait que :
    \[\fonction {\psi} {\mathscr C^0(\R,\R)} {\mathscr L(S,\R)} f {\varphi_f}\]

    est injective.
    À chaque fonction continue on associe donc une forme linéaire.
    Ceci généralise la notion de fonctions et définit la notion de \emph{distributions}.
\end{remarque}

\subsection{Complément hors programme sur les courbes paramétrées}

\subsubsection{Vocabulaire}

Soit $f : I \to E$.
\begin{itemize}
    \item Si $E = \R^2$, $f$ est appelée courbe plane
    \item Si $E = \R^3$, $f$ est appelée courbe gauche
\end{itemize}

\subsubsection{Interprétation cinématique}

Soit $f \in \mathscr C^2(I, E)$. 
Pour $t \in I$, $f'(t)$ est appelé vecteur vitesse à l'instant $t$ et $f''(t)$ est appelé vecteur accélération à l'instant $t$.

\subsubsection{Étude locale}

Soit $f \in \mathscr C^k(I, E), \, k \geqslant 2$ et $a \in I$. 

On suppose qu'il existe $p \in \ninter 1 k$ minimal tel que $f^{(p)}(a) \neq 0$.

De plus, on suppose qu'il existe $q \in \ninter p k$ minimal tel que $\left( f^{(p)}(a), f^{(q)}(a) \right)$ 
est libre.

\begin{itemize}
    \item Si $p = 1$ on dit que le point $f(a)$ est régulier. 
    \item Si $p = 1$ et $q = 2$, on dit que le point $f(a)$ est birégulier. 
\end{itemize}

\idee[-0.25cm]{Tout cela se démontrerait "facilement" avec des DL.}

La droite tangeante à la courbe est dirigée par $f^{(p)}(a)$ et 
il existe quatre types de points : 
\begin{itemize}
    \item $p$ impair, $q$ pair : point ordinaire, toute droite passant par $f(a)$ traverse alors la courbe, sauf la tangeante.
    \item $p$ impair, $q$ impair : point d'inflexion, toute droite passant par $f(a)$ traverse alors la courbe. 
    \item $p$ pair, $q$ impair : point de rembroussement de première espèce. 
    \item $p$ pair, $q$ pair : point de rembroussement de seconde espèce, aucune droite ne traverse la courbe.
\end{itemize}

\begin{figure}[htbp]
    \centering
    \begin{minipage}{0.45\textwidth}
        \centering
        \begin{tikzpicture}
            \begin{axis}[
                width=\linewidth,
                height=4.8cm,
                axis lines=middle,
                axis line style={-stealth, black!50, line width=0.4pt},
                title={Point ordinaire},
                xlabel={$x$},
                ylabel={$y$},
                label style={font=\footnotesize},
                xmin = -1.0, xmax = 1.0,
                ymin = -1.0, ymax = 1.0,
                ticks=none,
                samples=100,
            ]
                \addplot[black!50, line width=0.5pt, thick] {1.2*x};
                \addplot[vorange, thick] {x^2};

                \draw[->, thick, forestgreen] (axis cs:0,0) -- (axis cs:0.6,0)
                    node[below, text=forestgreen] {$f^{(p)}(t_0)$};

                \draw[->, thick, dangerred] (axis cs:0,0) -- (axis cs:0.35,{1.2*0.35})
                    node[left, xshift=-6pt, text=dangerred] {$f^{(q)}(t_0)$};
            \end{axis}
        \end{tikzpicture}
    \end{minipage}
    \hfill
    \begin{minipage}{0.45\textwidth}
        \centering
        \begin{tikzpicture}
            \begin{axis}[
                width=\linewidth,
                height=4.8cm,
                axis lines=middle,
                axis line style={-stealth, black!50, line width=0.4pt},
                title={Point d'inflexion},
                xlabel={$x$},
                ylabel={$y$},
                label style={font=\footnotesize},
                xmin = -1.0, xmax = 1.0,
                ymin = -1.0, ymax = 1.0,
                ticks=none,
                samples=100,
            ]
                \addplot[black!50, line width=0.5pt, thick] {1.2*x};
                \addplot[vorange, thick] {x^3};

                \draw[->, thick, forestgreen] (axis cs:0,0) -- (axis cs:0.6,0)
                    node[below, text=forestgreen] {$f^{(p)}(t_0)$};

                \draw[->, thick, dangerred] (axis cs:0,0) -- (axis cs:0.35,{1.2*0.35})
                    node[left, xshift=-6pt, text=dangerred] {$f^{(q)}(t_0)$};
            \end{axis}
        \end{tikzpicture}
    \end{minipage}
    \vfill
    \begin{minipage}{0.45\textwidth}
        \centering
        \begin{tikzpicture}
            \begin{axis}[
                width=\linewidth,
                height=4.8cm,
                axis lines=middle,
                axis line style={-stealth, black!50, line width=0.4pt},
                title={Rebroussement 1ère espèce},
                xlabel={$x$},
                ylabel={$y$},
                label style={font=\footnotesize},
                xmin = -1.0, xmax = 1.0,
                ymin = -1.0, ymax = 1.0,
                ticks=none,
                samples=100,
            ]
                \addplot[black!50, line width=0.5pt, thick] {1.0*x};

                \addplot[
                    domain=-0.95:0.95,
                    vorange,
                    very thick,
                    variable=\t
                ] ({t^2}, {0.9*t^3});


                \draw[->, thick, forestgreen] (axis cs:0,0) -- (axis cs:0.6,0)
                    node[below, text=forestgreen] {$f^{(p)}(t_0)$};

                \draw[->, thick, dangerred] (axis cs:0,0) -- (axis cs:0.35,{1.0*0.35})
                    node[left, xshift=-6pt, text=dangerred] {$f^{(q)}(t_0)$};
            \end{axis}
        \end{tikzpicture}
    \end{minipage}
    \hfill
    \begin{minipage}{0.45\textwidth}
        \centering
        \begin{tikzpicture}
            \begin{axis}[
                width=\linewidth,
                height=4.8cm,
                axis lines=middle,
                axis line style={-stealth, black!50, line width=0.4pt},
                title={Rebroussement 2ème espèce},
                xlabel={$x$},
                ylabel={$y$},
                label style={font=\footnotesize},
                xmin = -1.0, xmax = 1.0,
                ymin = -1.0, ymax = 1.0,
                ticks=none,
                samples=100,
            ]
                \addplot[black!50, line width=0.5pt, thick] {1.3*x};

                \addplot[
                    domain=-0.95:0.95,
                    vorange,
                    very thick,
                    variable=\t
                ] ({t^2 + 0.15*(t^4 + 0.35*t^5)}, {1.15*(t^4 + 0.35*t^5)});

                \draw[->, thick, forestgreen] (axis cs:0,0) -- (axis cs:0.6,0)
                    node[below, text=forestgreen] {$f^{(p)}(t_0)$};

                \draw[->, thick, dangerred] (axis cs:0,0) -- (axis cs:0.35,{1.3*0.35})
                    node[left, xshift=-6pt, text=dangerred] {$f^{(q)}(t_0)$};
                
            \end{axis}
        \end{tikzpicture}
    \end{minipage}
\end{figure}

\subsubsection{Étude globale d'une courbe plane}

$E = \R^2$, $\fonction f I {\R^2} t {\begin{pmatrix} x(t) \\ y(t) \end{pmatrix}}$ $f$ de classe $\mathscr C^k, \, k \geqslant 1$.

On trace le tableau 

$$\begin{array}{c|c}
    t & \cdots \\
    \hline
    x'(t) & \cdots \\
    \hline
    y'(t) & \cdots \\
    \hline
    y(t) & \cdots \\    
\end{array}$$

\section{Intégration sur un segment}
\subsection{Théorème-définition}

\begin{theoreme}{Intégrale de $f \in \mathscr C^{pm}$ à valeurs dans un evn de dim finie}{}
    Soit $f : [a, b] \to E$ et $B = (e_1, \ldots, e_n)$ une base de $E$.

    Notons $\displaystyle \forall t \in [a, b], \, f(t) = \sum_{i=1}^n f_i(t)e_i$.

    On suppose que $\forall i \in \ninter{1}{n}$, $f_i$ est continue par morceaux sur $[a, b]$.

    Cela permet de définir pour $i \in \ninter{1}{n} $, $\displaystyle \int_a^b f_i(t) \dt \in \K$.

    Alors si $B' = (e'_1, \ldots, e'_n)$ est une autre base de $E$, et si 
    $\displaystyle \forall t \in [a, b]$,  
    
    $\displaystyle f(t) = \sum_{i=1}^n g_i(t)e'_i$ avec 
    $g_i$ également continues par morceaux, on a :
    $$\sum_{i=1}^n \left( \int_a^b f_i(t) \dt \right) e_i = \sum_{i=1}^n \left( \int_a^b g_i(t) \dt \right) e'_i$$

    Ce vecteur est noté $\displaystyle \int_a^b f(t) \dt$.

    On dit que $f$ est continue par morceaux sur $[a, b]$. Cette notion ne dépend donc pas 
    du choix de la base.
\end{theoreme}

\idee{"Tout bête" (Quibel)

Sinon exprimer une base dans l'autre.}

\begin{proof}
    Soit $i \in \ninter 1 n$. 

    Décomposons $e_i = \displaystyle \sum_{j=1}^n a_{i,j} e'_j$.

    On a donc $\forall t \in [a, b]$, 

    \begin{align*}
        f(t) & = \sum_{i=1}^n f_i(t) \sum_{j=1}^n a_{i,j} e'_j \\
        & = \sum_{j=1}^n \left( \sum_{i=1}^n a_{i,j} f_i(t) \right) e'_j
    \end{align*}

    Par unicité des coordonnée $\displaystyle g_j(t) = \sum_{i=1}^n a_{i,j} f_i(t)$.

    Donc $g_j$ est continue par morceaux et 

    \begin{align*}
        \sum_{i=1}^n \left( \int_a^b g_i(t) \dt \right) e'_i & = \sum_{j=1}^n \left( \int_a^b \sum_{i=1}^n a_{i,j} f_i(t) \dt \right) e'_j \\
            & = \sum_{j=1}^n \int_a^b \sum_{i=1}^n a_{i,j} f_i(t) \dt e'_j \\
            & = \sum_{j=1}^n \sum_{i=1}^n a_{i,j} \left( \int_a^b f_i(t) \dt \right) e'_j \\
            & = \sum_{i=1}^n \left( \int_a^b f_i(t) \dt \right) \underbrace{\sum_{j=1}^n a_{i,j} e'_j}_{e_i} \\ 
    \end{align*}
\end{proof}

\begin{remarque}
Soit $f$ continue par morceaux sur $[a, b]$. Alors on note indifféremment :
$$\int_a^b f(t) \dt = \int_a^b f = \int_{[a, b]} f \in E$$
\end{remarque}

\subsection{Propriétés}

\begin{proposition}{Linéarité de l'intégrale}{}
    L'ensemble $\mathscr C^{pm}([a, b], E)$ est un $\K$-espace vectoriel et 
    $$ \fonction{\varphi}{\mathscr C^{pm}([a, b], E)}{E}{f}{\int_a^b f(t) \dt}$$ 
    est linéaire. 
\end{proposition}

\idee{Coordonnée par coordonnée.}

\begin{proof}
    Soient $f,g \in \mathscr C^{pm}([a,b],E)$, $\alpha,\beta \in \K$ et $B = (e_1, \ldots, e_n)$ une base de $E$, notons :
    
    $\displaystyle \forall t \in [a, b], \, f(t) = \sum_{i=1}^n f_i(t)e_i$ et
    $\displaystyle g(t) = \sum_{i=1}^n g_i(t)e_i$
    
    Alors :
    $\displaystyle \forall t \in [a, b], \, (\alpha f + \beta g)(t) = \sum_{i=1}^n (\alpha f_i(t) + \beta g_i(t))e_i$.

    Donc $\alpha f + \beta g$ est continue par morceaux (car on a la proposition dans $\K$) et par définition,

    \begin{align*}
        \int_a^b (\alpha f + \beta g)(t) \dt & = \sum_{i=1}^n \left( \int_a^b (\alpha f_i(t) + \beta g_i(t)) \dt \right) e_i\\
        & = \sum_{i=1}^n \left( \alpha \int_a^b f_i(t) \dt + \beta \int_a^b g_i(t) \dt \right) e_i\\
        & = \alpha \sum_{i=1}^n \left( \int_a^b f_i(t) \dt \right) e_i + \beta \sum_{i=1}^n \left( \int_a^b g_i(t) \dt \right) e_i\\
        & = \alpha \int_a^b f(t) \dt + \beta \int_a^b g(t) \dt
    \end{align*}
\end{proof}

\begin{proposition}{Notations}{}
    Soit $f : i \to E$ continue par morceaux sur tout segment de $I$. Soient $a, b \in I$.
    \begin{itemize}
        \item Si $a < b$, alors $\displaystyle \int_a^b f(t) \dt = \int_{[a, b]} f$
        \item Si $a = b$, par convention $\displaystyle \int_a^b f(t) \dt = 0_E$
        \item Si $a > b$, alors $\displaystyle \int_a^b f(t) \dt = - \int_{[b, a]} f$
    \end{itemize}
\end{proposition}


\begin{proposition}{Relation de Chasles}{}
    Soit $f : I \to E$ continue par morceaux sur tout segment de $I$. Soient $a, b, c \in I$.

    Alors $\displaystyle \int_a^c f(t) \dt = \int_a^b f(t) \dt + \int_b^c f(t) \dt$.
    
\end{proposition}

\begin{proof}
    Coordonnée par coordonnée.
\end{proof}

\marginenote{Remarque}{$f$ est dite continue par morceaux sur $I$ si elle l'est sur tout segment de $I$.}

\begin{theoreme}{Composition de l'intégrale par une application linéaire}{}
    Soit $f \in \mathscr C^{pm}([a, b], E)$, et soit $F$ un $\K$-espace vectoriel normé de dim finie, et 
    $L \in \mathscr L(E, F)$.

    Alors $L \circ f \in \mathscr C^{pm}([a, b], F)$ et $\displaystyle \int_a^b L(f(t)) \dt = L\left( \int_a^b f(t) \dt \right)$.
\end{theoreme}

\idee{Coordonnée par coordonnée, \emph{as usual}.}

\begin{proof}
    Soit $(e_1, \ldots, e_n)$ une base de $E$, et $(b_1, \ldots, b_q)$ une base de $F$.

    Posons $\displaystyle \forall t \in [a, b], \, f(t) = \sum_{i=1}^n f_i(t)e_i$.

    Donc $\displaystyle (L \circ f)(t) = \sum_{i=1}^n f_i(t)L(e_i)$.

    Décomposons pour tout $i$, $\displaystyle L(e_i) = \sum_{j=1}^q a_{j, i} b_j$.
    On a alors :

    \begin{align*}
        (L \circ f) (t) & = \sum_{i=1}^n f_i(t) \sum_{j=1}^q a_{j, i} b_j \\
        & = \sum_{j=1}^q \underbrace{\left( \sum_{i=1}^n a_{j, i} f_i(t) \right)}_{g_j(t)} b_j
    \end{align*}

    Les $f_i$ sont continues par morceaux donc $g_k$ aussi, donc $L \circ f$ est continue par morceaux et par définition : 
    {\allowdisplaybreaks
    \begin{align*}
        \int_a^b L \circ f & = \sum_{j=1}^q \left( \int_a^b g_j(t) \dt \right) b_j \\
        & = \sum_{j=1}^q \left( \int_a^b \sum_{i=1}^n a_{j, i} f_i(t) \dt \right) b_j \\
        & = \sum_{i = 1}^n \sum_{k = 1 }^{q } \int_{a }^{b } a_{k, i} f_i(t) \dt b_k \\
        & = \sum_{i=1}^n \left( \int_a^b f_i(t) \dt \right) \sum_{k=1}^q a_{k, i} b_k \\
        & = L \left(\left(\sum_{i = 1 }^{n } \int_{a }^{b } f_i(t) \dt\right) e_i\right) \\
        & = L \left( \int_a^b f(t) dt \right)
    \end{align*}
    }
\end{proof}

\subsection{Sommes de Riemann}

\begin{definition}{Somme de Riemann}{}
    Soit $N \geqslant 2$ et $f \in \mathscr C^{pm}([a, b], E)$.

    On appelle somme de Riemann à gauche (resp. à droite) pour $f$ sur $[a, b]$ associée à la 
    subdivision régulière en $N$ segments le vecteur :
    $$R_{g, N, f, [a, b]} = \frac{b-a} N \sum_{k=0}^{N-1} f\left( a + k \frac{b-a}{N} \right)$$

    Resp. $\displaystyle R_{d, N, f, [a, b]} = \frac{b-a} N \sum_{k=1}^{N} f\left( a + k \frac{b-a}{N} \right)$
\end{definition}

\begin{remarque}
    La notation avec les indices n'est pas standard.
    
    On peut aussi définir ces sommes sur des subdivisions non régulières et prendre un point quelconque dans les intervalles mais ce n'est pas utile ici.
\end{remarque}

\begin{proposition}{}{}
    Soit $f \in \mathscr C^{pm}([a, b], E)$. Alors les sommes de Riemann pour $f$ convergent vers 
    $\displaystyle \int_a^b f(t) \dt$ quand $N \to +\infty$.

    On a donc :
    $$ \frac{b - a }{N } \sum_{k = 0 }^{N - 1 } f\left( a + k \frac{b - a } N \right) \limit{N}{+ \infty} \int_a^b f(t) \dt$$
    $$ \frac{b - a }{N } \sum_{k = 1 }^{N} f\left( a + k \frac{b - a } N \right) \limit{N}{+ \infty} \int_a^b f(t) \dt$$

\end{proposition}

\idee{Surprise surprise, coordonnée par coordonnée.}

\begin{proof}
Soit $B= (e_1,\dots,e_n)$ une base de $E$.
    \begin{align*}
        R_{g, N, f, [a, b]} & = \frac{b - a }{N } \sum_{k = 0 }^{N - 1 } \sum_{i = 1 }^{n } f_i\left( a + k \frac{b - a } N \right) e_i \\
        & = \sum_{i = 1 }^{n } \left( \frac{b - a }{N } \sum_{k = 0 }^{N - 1 } f_i\left( a + k \frac{b - a } N \right) \right) e_i \\
        & \limit{N}{+ \infty} \sum_{i = 1 }^{n } \left( \int_a^b f_i(t) \dt \right) e_i \\
        & = \int_a^b f(t) \dt
    \end{align*}
De même pour la somme à droite.
\end{proof}

\begin{theoreme}{Inégalité triangulaire pour l'intégrale}{}
    Munissons $E$ d'une norme $||\cdot||$. Soit $f \in \mathscr C^{pm}([a, b], E)$.
    Alors $||f|| \in \mathscr C^{pm}([a, b], \R)$ et :
    $$\left|\left| \int_a^b f(t) \dt \right|\right| \leqslant \int_a^b ||f(t)|| dt$$
\end{theoreme}

\idee{Sommes de Riemann et passage à la limite.}

\begin{proof}
    $||f||$ est continue par morceaux par composition car $||\cdot||$ est continue
    et $f$ est continue par morceaux (prendre une subdivision adaptée pour le prouver).

    Soit $N \geqslant 2$, alors,

    \begin{align*}
        \left|\left| R_{g, N, f, [a, b]} \right|\right| & = \left|\left| \frac{b-a}{N} \sum_{k=0}^{N-1} f\left( a + k \frac{b-a}{N} \right) \right|\right| \\
        & \leqslant \frac{b-a}{N} \sum_{k=0}^{N-1} \left|\left| f\left( a + k \frac{b-a}{N} \right) \right|\right| \\
        & \leqslant R_{g, N, ||f||, [a, b]}
    \end{align*}

    Et donc par passage à la limite, 
    $$\left|\left| \int_a^b f(t) \dt \right|\right| \leqslant \int_a^b ||f(t)|| \dt$$
\end{proof}


\subsection{Intégrale de la borne supérieure}

\begin{theoreme}{Primitive fonction vectorielle}{}
    Soit $f \in \mathscr C(I, E)$ et $a \in I$. 

    Alors $\fonction{F}{I}{E}{x}{\int_a^x f(t) \dt}$ est de classe $\mathscr C^1$ sur $I$ et 
    vérifie $F' = f$.

    $F$ est appelée la primitive de $f$ sur $I$ qui s'annule en $a$. 
\end{theoreme}

\idee{Coordonnée par coordonnée.}

\begin{proof}
    Notons $B = (e_1, \ldots, e_n)$ une base de $E$, et $\forall t \in I, \, f(t) \displaystyle = \sum_{i = 0 }^{ n } f_i(t) e_i$

    Donc $\displaystyle F(x) = \sum_{i = 0 }^{ n } \left( \int_a^x f_i(t) \dt \right) e_i$.

    Par le théorème fondamental pour les fonctions réelles, $F_i$ est de classe $\mathscr C^1$ et $F_i' = f_i$.
\end{proof}

\begin{corollaire}{Calcul d'intégrales}{}
    Soit $f \in \mathscr C(I, E)$ et $F \in \mathscr C^1(I, E)$ tel que 
    $F' = f$, alors $F$ est appelée une primitive de $f$. Alors $\forall a, b \in I$, 
    
    $$\displaystyle \int_a^b f(t) \dt = F(b) - F(a)$$
\end{corollaire}

\idee{Prendre deux primitives différentes.}

\begin{proof}
    Posons $\ds \forall x \in I, \, G(x) = \int_a^x f(t) dt$.

    Par le théorème, $G' = f = F'$ donc $(F - G)' = 0$ et donc $F - G$ est constante, 
    donc $\ds F(b) - F(a) = G(b) - G(a) = \int_a^b f(t) \dt$.
\end{proof}

\subsubsection{Recherche de primitives}

$$\renewcommand{\arraystretch}{1.6}
\begin{array}{|c|c|c|}
    \hline
    I & f(x) & \int_{a }^{x } f(t) \dt \\
    \hline
    \R & x^n & \frac{x^{n+1}}{n + 1} \\
    \hline
    \R^*_+ & x^\alpha, \alpha \neq -1 & \frac{x^{\alpha + 1}}{\alpha + 1} \\
    \hline
    \R^*_+ \text{ ou } \R_-^*& \frac 1 x & \ln |x| \\
    \hline
    \R & e^x & e^x \\
    \hline
    \R & \cos x & \sin x \\
    \hline
    \R & \sin x & -\cos x \\
    \hline
    \R & \cosh x & \sinh x \\
    \hline
    \R & \sinh x & \cosh x \\
    \hline
    \R & \frac 1 {1 + x^2} & \arctan x \\
    \hline
    ]-1, 1[ & \frac 1 {\sqrt{1 - x^2}} & \arcsin x\\
    \hline
\end{array}$$

\subsubsection{Intégration par parties}

Soient $u, v \in \mathscr C^1([a, b], \K)$.
Alors :
$$\int_{a }^{b } u'\cdot v = \left[ u \cdot v \right]_{a }^{b } - \int_{a }^{b } u \cdot v'$$

\subsubsection{Changement de variable}

Soit $\varphi \in \mathscr C^1(J, I)$ et $\alpha, \beta \in J$, 
$\varphi(\alpha) = a, \, \varphi(\beta) = b$.

Soit $f \in \mathscr C(I, E)$. Alors $$\int_\alpha^\beta f(\varphi(u)) \varphi'(u) \dif u = \int_a^b f(t) \dt$$
En posant $t = \varphi(u)$, donc $\dx = \varphi'(u) \dif u$.

Changements de variable usuels : 

$$\renewcommand{\arraystretch}{2.1}
{\everymath{\displaystyle}
\begin{array}{|c|c|}
    \hline 
    \text{\textbf{Intégrale}} & \text{\textbf{changement}} \\
    \hline
    \displaystyle \int_a^b f(t^2)t \dt & x = t^2 \; \dx = 2t \dt \\
    \hline
    \int_a^b f(\cos t) \sin t \dt & x = \cos t \; \dx = -\sin t \dt \\
    \hline
    \int_a^b F(e^t) \dt & \substack{x = e^t \; \dx = e^t \dt \\ t = \ln(x) \; \dt = \frac{\dx} x} \\
    \hline 
    \int_{a }^{b } F(\cos t, \sin t) \dt, \; F \in \K(X) & \substack{u = \tan \frac t 2, \; t = 2 \arctan u \; \dt = \frac{2 \dif u}{1 + u^2} \\ \cos t = \frac{1 - u^2}{1 + u^2} \; \sin t = \frac{2u}{1 + u^2}} \\
    \hline
    \int_{a }^{b } F(x, \sqrt{1 - x^2}) \dx & x = \sin t \\
    \hline
    \int_{a }^{b } F(x, \sqrt{x^2 - 1}) \dx & x = \cosh t \\
    \hline
    \int_{a }^{b } F(x, \sqrt{1 + x^2}) \dx & x = \sinh t \\
    \hline
\end{array}}$$

\subsubsection{Décomposition en éléments simples}

\begin{exemple}
    $\displaystyle I = \int_0^x \sqrt{\tan t} \text d t$, $x \in [0, \frac \pi 2[$.

    Posons $u = \sqrt{\tan t}$, donc $t = \arctan(u^2)$, $\dt = \frac{2u}{1 + u^4} \dif u$.

    Donc $\displaystyle I = \int_0^{\sqrt{\tan x}} \frac{2u^2}{1 + u^4} \dif u$.

    Décomposition en éléments simples : $\frac{2X^2}{X^4 + 1}$

    $1 + X^4 = (1 + X^2)^2 -2X^2 = (1 - \sqrt 2 X + X^2)(1 + \sqrt 2 X + X^2)$

    Donc $\displaystyle \frac{2X^2}{1 + X^4} = \frac{aX + b}{X^2 - \sqrt 2 X + 1} + \frac{cX + d}{X^2 + \sqrt 2 X + 1}$.

    En évaluant en 0 : $b + d = 0$.

    Limite en $+\infty : 0 = a + c$. 

    On a donc : $\displaystyle \frac{2X^2}{1 + X^4} = \frac{aX + b}{X^2 - \sqrt 2 X + 1} - \frac{aX + b}{X^2 + \sqrt 2 X + 1} = \frac{(aX + b)(2 \sqrt 2 X)}{X^4 + 1}$.

    Donc $b = 0$ et $a = \frac{\sqrt 2 }{2}$. Donc $\forall u, $

    \begin{align*}
        \frac{2 u^2}{1 + u^4} & = \frac 1 2 \left[\frac{\sqrt 2 u }{u^2 - \sqrt 2 u + 1} - \frac{\sqrt 2 u }{u^2 + \sqrt 2 u + 1}\right]\\
        & = \frac 1 2 \left[ \frac 1 {\sqrt 2} \frac{2u - \sqrt 2}{u^2 - \sqrt 2 u + 1} + \frac 1 {u^2 - \sqrt 2 u + 1} - \frac 1 {\sqrt 2} \frac{2u + \sqrt 2}{u^2 + \sqrt 2 u + 1} + \frac 1 {u^2 + \sqrt 2 u + 1} \right]\\
        & = \frac 1 {2 \sqrt 2} \left[ \frac{2u - \sqrt 2}{u^2 - \sqrt 2 u + 1} - \frac{2u + \sqrt 2}{u^2 + \sqrt 2 u + 1} \right] + \frac 1 2 \left[ \frac{1}{\left( u - \frac{\sqrt 2 }{2} \right)^2 + \frac 1 2} + \frac{1}{\left( u + \frac{\sqrt 2}{2} \right)^2 + \frac 1 2} \right]\\
        & = \frac 1 {2 \sqrt 2} \left[ \frac{2u - \sqrt 2}{u^2 - \sqrt 2 u + 1} + \frac 1 {(\sqrt 2 u - 1)^2 + 1} + \frac 1 {(\sqrt 2 u + 1)^2 + 1} \right]
    \end{align*}

    Donc $\displaystyle I = \left[ \frac{1 }{2 \sqrt 2 } \ln(u^2 - \sqrt 2 u + 1) - \ln(u^2 + \sqrt 2 u + 1) + \frac{\arctan(\sqrt 2 u - 1)}{\sqrt 2} + \frac{\arctan(\sqrt 2 u + 1)}{\sqrt 2} \right]_0^{\sqrt{\tan x}}$
\end{exemple}

\begin{theoreme}{Inégalité des accroissements finis}{}
    Soit $f \in \mathscr C^1(I, E)$, soient $\alpha,\beta \in I$. Soit $M > 0, \, \forall x \in [\alpha,\beta], \, ||f'(x)|| \leqslant M$.

    Alors $\forall a, b \in [\alpha,\beta]$, 

    $$||f(b) - f(a)|| \leqslant M |b - a|$$
\end{theoreme}

\idee{Inégalité triangulaire des intégrales.}

\begin{proof}

    $\ds f(b) - f(a) = \int_a^b f'(t) \dt$.

    Donc $\displaystyle ||f(b) - f(a)|| \leqslant \int_a^b ||f'(t)|| \dt \leqslant \int_a^b M \dt = M |b - a|$.
\end{proof}

\begin{remarque}
    À contrario, le théorème de Rolle et le théorème des accroissement finis ne sont pas valables pour les fonctions vectorielles.
\end{remarque}

\subsection{Formules de Taylor}

\marginenote{Histoire !}{À l'époque, Taylor notait les dérivées successives par des lettres $m,n,\dots$ car la notation $f^{(n)}$ n'existait pas.}

\begin{proposition}{Formule de Taylor-Young}{}
    Soit $f \in \mathscr C^n(I, E)$, alors pour $a \in I$, $h \in \R$ tel que $a + h \in I$, on a :
    $$f(a + h) = \sum_{k = 0 }^{n } \frac{h^k}{k!} f^{(k)}(a) + o(h^n)$$
\end{proposition}

\begin{proof}
    Composante par composante.
\end{proof}

\begin{proposition}{Formule de Taylor avec reste intégrale}{}
    Soit $f \in \mathscr C^{n + 1}(I, E)$ et $a, b \in I$. 

    $$f(b) = \sum_{k = 0 }^{n } \frac{(b - a)^k}{k!} f^{(k)}(a) + \int_a^b \frac{(b - t)^n}{n!} f^{(n+1)}(t) dt$$
\end{proposition}

\begin{proof}
    Composante par composante.
\end{proof}

\begin{proposition}{Inégalité de Taylor-Lagrange}{}
    Soit $f \in \mathscr C^{n + 1}(I, E)$, $a, b \in I$.  Soit $M > 0, \, \forall t \in [a, b], \, ||f^{(n + 1)}(t)|| \leqslant M$.
    Alors :
    $$\left|\left| f(b) - \sum_{k = 0 }^{n } \frac{(b - a)^k}{k!} f^{(k)}(a) \right|\right| \leqslant \frac{M |b - a|^{n + 1}}{(n + 1)!}$$
\end{proposition}

\begin{proof}
    On prend $a < b$.

    \begin{align*}
    \left|\left| f(b) - \sum_{k = 0}^n \frac{(b - a)^k}{k!} f^{(k)}(a) \right|\right|  & = \left|\left| \int_a^b \frac{f^{(n + 1)}(t)}{n!} (b - t)^n \dt \right|\right| \\
     & \leqslant \int_a^b \left| \frac{(b - t)^n}{n!} \right| M \dt \\
        & \leqslant M \int_{a }^{b } \frac{(b - t)^n}{n!} \dt \\
        & = M \frac{(b - a )^{n + 1}}{(n+ 1)!}
    \end{align*}
\end{proof}
